\title{The Era of 1-bit LLMs: All Large Language Models are in 1.58 Bits}

\begin{abstract}
Recent advances, exemplified by BitNet~\cite{bitnet}, are ushering in a new generation of 1-bit Large Language Models (LLMs). In this paper, we present a 1-bit LLM variant, called \textbf{\text{BitNet b1.58}}, where each parameter (or weight) of the LLM is ternarized to \{-1, 0, 1\}. This model achieves comparable perplexity and downstream task performance to the full-precision (FP16 or BF16) Transformer LLM with equivalent model size and training data, while offering substantial improvements in latency, memory usage, throughput, and energy efficiency. More importantly, the 1.58-bit LLM establishes a novel scaling law and training methodology for developing future LLMs that are both efficient and high-performing. Additionally, it introduces a new computational paradigm and paves the way for designing dedicated hardware tailored to 1-bit LLMs.
\end{abstract}

\section{The Era of 1-bit LLMs}

The AI community has recently witnessed rapid advancements in the scale and capabilities of Large Language Models (LLMs). These models have achieved impressive results across various natural language processing tasks, yet their growing size introduces deployment difficulties and raises environmental and financial concerns due to elevated energy demands.

A common strategy to mitigate these issues involves applying post-training quantization to create low-bit models for inference~\cite{smoothquant, gptq, quip, quip_sharp}. This method lowers the precision of weights and activations, markedly decreasing the memory footprint and computational load of LLMs. The trend has shifted from 16-bit precision toward lower-bit formats such as 4-bit models~\cite{gptq, awq}. Nonetheless, post-training quantization remains suboptimal despite its widespread adoption in industrial LLM implementations.

Recent research on 1-bit model architectures, such as BitNet~\cite{bitnet}, offers a promising avenue for lowering the costs associated with LLMs while preserving their effectiveness. Standard LLMs operate using 16-bit floating-point values (i.e., FP16 or BF16), with the majority of their computation dominated by matrix multiplication. Consequently, the primary computational expense arises from floating-point addition and multiplication operations. In contrast, BitNet’s matrix multiplication involves only integer addition, resulting in substantial reductions in energy consumption for LLMs. Since power is often the fundamental constraint on compute performance in many chips, these energy savings can also translate into faster processing speeds.

Beyond computation, the transfer of model parameters from DRAM to on-chip accelerator memory (e.g., SRAM) can be costly during inference. Attempts to enlarge SRAM to enhance throughput have been made, but these come with significantly higher costs compared to DRAM. Relative to full-precision models, 1-bit LLMs drastically reduce memory requirements in both capacity and bandwidth. This reduction can considerably lower the cost and duration of loading weights from DRAM, thereby enabling faster and more efficient inference.

In this paper, we propose an important 1-bit LLM variant named \textbf{\text{BitNet b1.58}}, where each parameter is ternary and takes values from the set \{-1, 0, 1\}. By introducing an additional value 0 to the original 1-bit BitNet, we achieve an effective 1.58-bit representation in binary. \text{BitNet b1.58} preserves all advantages of the original 1-bit BitNet, including a novel computation paradigm that almost entirely eliminates multiplication operations during matrix multiplication and is highly optimizable. Moreover, it maintains the same energy efficiency as the original 1-bit BitNet while being significantly more efficient in terms of memory usage, throughput, and latency compared to FP16 LLM baselines. In addition, \text{BitNet b1.58} presents two further benefits. First, its modeling power is enhanced through explicit feature filtering, enabled by the inclusion of zero in the model weights, which can substantially boost 1-bit LLM performance. Second, our experiments demonstrate that \text{BitNet b1.58} can match full-precision (i.e., FP16) baselines in both perplexity and end-task metrics from a 3B model scale onwards, given the same settings (e.g., model size, training tokens, etc.).

\section{BitNet b1.58}

\text{BitNet b1.58} builds upon the BitNet architecture, a Transformer model that substitutes \emph{nn.Linear} layers with \emph{BitLinear} layers. It is trained from scratch, employing 1.58-bit precision for weights and 8-bit precision for activations. In comparison to the original BitNet, it introduces several changes, which we outline below.

\paragraph{Quantization Function.} To restrict the weights to the set \{-1, 0, +1\}, we utilize an \emph{absmean} quantization function. This function initially normalizes the weight matrix by its mean absolute value, then rounds each element to the closest integer within \{-1, 0, +1\}:

\begin{equation}
    \widetilde{W} = \mathrm{RoundClip}\left(\frac{W}{\gamma + \epsilon}, -1, 1\right),
\end{equation}

\begin{equation}
    \mathrm{RoundClip}(x, a, b) = \max\left(a, \min\left(b, \mathrm{round}(x)\right)\right),
\end{equation}

\begin{equation}
    \gamma = \frac{1}{nm}\sum_{ij} |W_{ij}|.
\end{equation}

The activation quantization function follows the same approach as in BitNet, with one key difference: instead of scaling activations prior to nonlinearities into the range $[0, Q_b]$, activations are scaled per token into $[-Q_b, Q_b]$, eliminating the need for zero-point quantization. This approach simplifies implementation and system-level optimization while having minimal impact on performance based on our experiments.

\paragraph{LLaMA-alike Components.}

The LLaMA architecture~\cite{llama, llama2} has become the standard backbone for open-source large language models. To align with the open-source ecosystem, \text{BitNet b1.58} incorporates components inspired by LLaMA. In particular, it employs RMSNorm~\cite{rmsnorm}, SwiGLU~\cite{swiglu}, rotary embeddings~\cite{rope}, and removes all bias terms. This design facilitates easy integration of \text{BitNet b1.58} into popular open-source frameworks such as Huggingface, vLLM~\cite{vllm}, and llama.cpp\footnote{\url{https://github.com/ggerganov/llama.cpp}} with minimal modifications.

\section{Results}

We conducted a comparison between \text{BitNet b1.58} and our re-implemented FP16 \text{LLaMA LLM} across multiple model sizes. To ensure a consistent basis for comparison, both models were pre-trained on the RedPajama dataset~\cite{redpajama} using 100 billion tokens. Their zero-shot capabilities were assessed across a variety of language benchmarks, including ARC-Easy~\cite{arc}, ARC-Challenge~\cite{arc}, Hellaswag~\cite{hellaswag}, Winogrande~\cite{winoGrande}, PIQA~\cite{piqa}, OpenbookQA~\cite{openbookqa}, and BoolQ~\cite{boolq}. Additionally, we recorded the validation perplexity on WikiText2~\cite{wikitext2} and C4~\cite{c4} datasets.

We also measured the GPU memory usage and latency during runtime for both \text{LLaMA LLM} and \text{BitNet b1.58}. These metrics were obtained using the FasterTransformer\footnote{\url{https://github.com/NVIDIA/FasterTransformer}} framework, which is optimized for reducing LLM inference latency on GPU hardware. The 2-bit kernel implementation from Ladder~\cite{ladder} was incorporated for \text{BitNet b1.58}. The primary metric reported was the time per output token, reflecting the dominant cost in inference.

Table~\ref{tab:ppl} presents a summary of the perplexity scores alongside resource costs for \text{BitNet b1.58} and \text{LLaMA LLM}. It indicates that \text{BitNet b1.58} reaches parity with the full-precision \text{LLaMA LLM} in perplexity starting at the 3B parameter model size, while being 2.71 times faster and requiring 3.55 times less GPU memory. Notably, the 3.9B \text{BitNet b1.58} model operates 2.4 times faster and uses 3.32 times less memory, yet delivers substantially improved performance compared to the \text{LLaMA LLM} 3B model.

Detailed zero-shot accuracy results on the downstream tasks are provided in Table~\ref{tab:zero-shot}. Our evaluation followed the procedure outlined in the \emph{lm-evaluation-harness}\footnote{\url{https://github.com/EleutherAI/lm-evaluation-harness}}. The findings demonstrate that the performance disparity between \text{BitNet b1.58} and \text{LLaMA LLM} diminishes with increasing model size. Crucially, \text{BitNet b1.58} matches the full-precision baseline starting at the 3B scale. Mirroring the perplexity trends, the task-specific results confirm that the 3.9B \text{BitNet b1.58} model surpasses the 3B \text{LLaMA LLM} while incurring lower memory and latency overhead. This validates that \text{BitNet b1.58} achieves a Pareto improvement over existing state-of-the-art LLMs.

\paragraph{Memory and Latency}

We additionally increased the model size to 7B, 13B, and 70B and assessed the associated costs. Figure~\ref{fig:latency-memory-energy} depicts the latency and memory usage trends, indicating that the speed-up improves as the model size grows. Notably, \text{BitNet b1.58} 70B achieves a speed that is 4.1 times faster than the \text{LLaMA LLM} baseline. This improvement arises because the time required for \emph{nn.Linear} operations expands with the model size. Memory usage exhibits a comparable pattern, since the embedding remains in full precision and its memory share decreases for larger models. Both latency and memory measurements were taken using a 2-bit kernel, suggesting potential for further optimization to reduce costs.

\paragraph{Energy}

We also estimated the energy consumption of arithmetic operations for both \text{BitNet b1.58} and \text{LLaMA LLM}. Our primary focus was on the matrix multiplication calculations, as these dominate the computational cost in LLMs. Figure~\ref{fig:energy} presents the breakdown of energy consumption. The bulk of \text{BitNet b1.58}'s computations involve INT8 additions, whereas \text{LLaMA LLM} operations include both FP16 additions and multiplications. Based on the energy model from~\cite{energycost, pokebnn}, \text{BitNet b1.58} achieves a 71.4-fold reduction in arithmetic operation energy consumption for matrix multiplication on 7nm chips. We also reported the overall energy expenditure for processing 512 tokens. Our findings reveal that as model size increases, \text{BitNet b1.58} becomes progressively more energy-efficient compared to the FP16 \text{LLaMA LLM} baseline. This efficiency gain is attributable to the growing proportion of \emph{nn.Linear} operations with model scale, while energy costs from other components diminish in larger models.

\paragraph{Throughput}

We evaluate the throughput of \text{BitNet b1.58} compared to \text{LLaMA LLM} with 70 billion parameters on two 80GB A100 GPUs, utilizing pipeline parallelism~\cite{gpipe} to enable running the \text{LLaMA LLM} 70B on these devices. The batch size was incrementally increased until the GPU memory was fully utilized, with a fixed sequence length of 512. As shown in Table~\ref{tab:throughput}, \text{BitNet b1.58} 70B can handle a batch size up to 11 times larger than \text{LLaMA LLM}, leading to a throughput improvement by a factor of 8.9.

\textbf{\text{BitNet b1.58} introduces a novel scaling law linking model performance and inference cost}. For context, the equivalences between model sizes in 1.58-bit and 16-bit precision, derived from Figures~\ref{fig:latency-memory-energy} and \ref{fig:energy}, are as follows:

\begin{itemize}
    \item 13B \text{BitNet b1.58} is more efficient in latency, memory consumption, and energy use compared to a 3B FP16 LLM.
    \item 30B \text{BitNet b1.58} surpasses a 7B FP16 LLM in terms of latency, memory, and energy efficiency.
    \item 70B \text{BitNet b1.58} outperforms a 13B FP16 LLM with respect to latency, memory, and energy consumption.
\end{itemize}

\paragraph{Training with 2T Tokens}

The total number of training tokens is a vital parameter for LLMs. To assess the token-scale scalability of \text{BitNet b1.58}, we trained a \text{BitNet b1.58} model on 2 trillion tokens, following the data pipeline of StableLM-3B~\cite{StableLM-3B-4E1T}, which represents the leading open-source 3B model. Both models were tested on a suite of benchmarks including Winogrande~\cite{winoGrande}, PIQA~\cite{piqa}, SciQ~\cite{sciq}, LAMBADA~\cite{lambada}, and ARC-easy~\cite{arc}. Zero-shot accuracy metrics are reported in Table~\ref{tab:2t}. For tasks evaluated via accuracy and normalized accuracy, we present the mean of the two. The StableLM 3B results at 2T tokens are sourced directly from its official technical report. Our results demonstrate that \text{BitNet b1.58} attains superior performance across all evaluated tasks, signifying that 1.58-bit LLMs also exhibit strong generalization capabilities.

\section{Discussion and Future Work}

\textbf{1-bit Mixture-of-Experts (MoE) LLMs}

Mixture-of-Experts (MoE) models have shown to be an efficient method for large language models (LLMs) by substantially lowering computational FLOPs. However, their deployment and practical use are hindered by high memory requirements and the overhead from inter-chip communication. These issues can be mitigated through the use of 1.58-bit LLMs. Specifically, the reduced memory footprint decreases the number of devices needed to host MoE models. Additionally, it significantly cuts down the cost of transferring activations over networks. Ideally, if the entire model fits on a single chip, this overhead would be eliminated altogether.

\textbf{Native Support of Long Sequence in LLMs}

With the advent of LLMs, managing long sequences natively has become increasingly important. A key obstacle for long sequence inference is the substantial memory used by KV caches. \text{BitNet b1.58} marks a notable advancement toward native long-sequence support by lowering activations from 16 bits to 8 bits, effectively enabling the doubling of context length within the same resource constraints. Furthermore, this representation can be compressed losslessly to 4 bits or even lower for 1.58-bit LLMs, a direction reserved for future investigation.

\textbf{LLMs on Edge and Mobile}

Applying 1.58-bit LLMs holds great promise for enhancing language model performance on edge and mobile platforms. These devices typically face limitations in memory capacity and computational power, restricting the size and effectiveness of LLMs they can support. The reduced memory usage and energy demand of 1.58-bit LLMs allow their deployment on such devices, unlocking numerous applications that were previously unattainable. This can significantly boost the functionality of edge and mobile devices and open new possibilities for LLM-based applications. Moreover, 1.58-bit LLMs are better suited for CPU architectures, which are prevalent in edge and mobile hardware, making \text{BitNet b1.58} highly efficient to run on these platforms, thereby enhancing their performance and capabilities.

\textbf{New Hardware for 1-bit LLMs}

Recent developments, such as those from~Groq\footnote{\url{https://groq.com/}}, have shown encouraging outcomes and substantial potential in developing specialized hardware (e.g., LPUs) tailored for LLMs. Building on this progress, we envision and advocate for the design of new hardware and systems specifically optimized for 1-bit LLMs, leveraging the novel computation paradigm introduced by BitNet~\cite{bitnet}. 

\end{document}